\documentclass[11pt]{article}
\usepackage[a4paper,margin=25mm]{geometry}
\usepackage[T1]{fontenc}
\usepackage{amsmath}
\usepackage{booktabs}
\usepackage{graphicx}
\usepackage{longtable}
\usepackage{listings}
\usepackage{xcolor}
\usepackage{hyperref}

\hypersetup{
  colorlinks=true,
  linkcolor=blue,
  urlcolor=blue,
  citecolor=blue
}
\lstset{
  basicstyle=\ttfamily\small,
  breaklines=true,
  columns=fullflexible,
  keepspaces=true,
  showstringspaces=false
}
\title{Git and GitHub for Physics\\One Shared Guide for Every Operating System}
\author{UCC Physics Society}
\begin{document}
\maketitle
\tableofcontents
\newpage


\section{Git is not GitHub}
Git records project history locally. GitHub hosts repositories online and
adds pull requests and reviews. Saving, committing, and pushing are different
operations; local Git works offline. The commands here work in PowerShell,
Bash, and the VS Code integrated terminal. For a command reference, see the
official \href{https://git-scm.com/docs}{Git documentation} and
\href{https://git-scm.com/book/en/v2}{Pro Git}.


\section{Identity and authentication}
Check \texttt{git --version}, then configure the identity used on new commits:
If Git is missing, follow the installation instructions in your
\href{../README.md}{operating-system guide} or use the
\href{https://git-scm.com/downloads}{official Git download page}.

\begin{lstlisting}
git config --global user.name "Your Name"
git config --global user.email "YOUR_CHOSEN_EMAIL"
\end{lstlisting}
Replace both placeholders. If you want email privacy, obtain your GitHub
no-reply address from account settings \emph{before} making public commits.
For one repository, omit \texttt{--global} and run inside it.
Never alter another user's global configuration on a shared login.


Author identity is \textbf{not login}:
\begin{description}
\item[HTTPS:] Use an approved credential manager/browser sign-in, or a narrowly
scoped token when required. An account password is not a Git HTTPS password.
\item[SSH:] Follow
\href{https://docs.github.com/en/authentication/connecting-to-github-with-ssh}{GitHub's SSH guide}.
Add only your public key to GitHub; keep the private key private.
\end{description}
Do not put tokens in remote URLs, notebooks, screenshots, or source.
Do not disable TLS or host-key checks to bypass authentication problems.

\section{Mental model and inspection commands}
\begin{center}
Working tree $\xrightarrow{\texttt{git add}}$ staging area
$\xrightarrow{\texttt{git commit}}$ local history
\end{center}
Pushing sends local commits to a remote repository; it does not automatically
include uncommitted edits.
\begin{longtable}{p{0.36\linewidth}p{0.57\linewidth}}
\toprule Command & Meaning\\ \midrule
\texttt{git status} & What changed, what is staged, current branch\\
\texttt{git diff} & Unstaged changes\\
\texttt{git diff --staged} & Exact changes proposed for the next commit\\
\texttt{git log --oneline -5} & Recent commits\\
\texttt{git remote -v} & Remote names and URLs; check for credentials before sharing\\
\bottomrule
\end{longtable}
Commits are deliberate snapshots, not automatic backups of every edit.

\section{Local practice lab}\label{sec:local-lab}
Work in a new writable directory \textbf{outside the workshop repository}:
\begin{lstlisting}
mkdir physics-git-practice
cd physics-git-practice
git init -b main
\end{lstlisting}
If an older Git does not support \texttt{-b}, use \texttt{git init} followed
by \texttt{git branch -M main}.
Create \texttt{README.md} in your editor:
\begin{lstlisting}
# Pendulum lab
We will measure the period of a pendulum.
\end{lstlisting}
Then run:
\begin{lstlisting}
git status
git add README.md
git diff --staged
git commit -m "Describe the pendulum experiment"
\end{lstlisting}
Add a second sentence about recording units, then:
\begin{lstlisting}
git diff
git add README.md
git diff --staged
git commit -m "Record the measurement units"
git log --oneline -5
git status
\end{lstlisting}
\textbf{Success:} two commits, a clean working tree, and an explanation of the
second change. If identity is missing, configure it as above and retry the commit.

\section{Ignore generated and private files}
Create \texttt{.gitignore} in your \emph{practice repository}:
\begin{lstlisting}
.venv/
__pycache__/
.ipynb_checkpoints/
.env
*.aux
*.log
*.out
*.toc
*.fls
*.fdb_latexmk
*.synctex.gz
*.blg
\end{lstlisting}
Do not blindly ignore all PDFs or CSVs. Selected release PDFs and small public
datasets can be intentional deliverables.

Ignoring affects untracked files only: it does not remove tracked files or
erase history. Do not treat it as secret protection.
Stage named files rather than immediately running \texttt{git add .};
review \texttt{git status} and \texttt{git diff --staged} for private data,
large output, and unrelated changes.

\section{Publish your own project}
Create an \textbf{empty} GitHub repository for this route: do not initialise it
with a README, licence, or ignore file. Replace the URL placeholders:
\begin{lstlisting}
git remote add origin https://github.com/YOUR_USER/YOUR_REPO.git
git remote -v
git push -u origin main
\end{lstlisting}
To work on an \emph{existing} project instead, run \texttt{git clone URL}
in its parent directory. Do not initialise inside an existing clone.
A ZIP download supplies files but no Git history.

\section{Branches and pull requests}
From a clean working tree in your practice repository:
\begin{lstlisting}
git switch -c improve-method
\end{lstlisting}
Edit the README, inspect the diff, then:
\begin{lstlisting}
git add README.md
git commit -m "Explain how the period is measured"
git push -u origin improve-method
\end{lstlisting}
On GitHub open a pull request into \texttt{main}. Explain purpose, changes,
and checks. Ask a partner to review; after merging:
\begin{lstlisting}
git switch main
git pull --ff-only
\end{lstlisting}
Without a remote, practise locally using \texttt{git switch main} and
\texttt{git merge improve-method} instead. Do not blindly combine both routes.
Update before beginning new shared work. Commit or deliberately stash unfinished
work before switching context. Never force-push to fix a beginner sync problem.

\section{Merge conflict practice}
Use only the disposable repository, starting clean:
\begin{enumerate}
\item Create \texttt{conflict-demo} with \texttt{git switch -c conflict-demo}.
\item Change an existing README sentence and commit it.
\item Switch to \texttt{main}; change that same sentence differently and commit.
\item Run \texttt{git merge conflict-demo}. Read the conflict message.
\item Replace the \texttt{<<<<<<<}, \texttt{=======}, and
\texttt{>>>>>>>} block with the intended final text.
\item Review the whole result, run \texttt{git add README.md}, then
\texttt{git commit}.
\item Check \texttt{git status} and
\texttt{git log --oneline --graph --all}.
\end{enumerate}
To cancel the unfinished practice merge, use \texttt{git merge --abort}.
Beginning clean makes this straightforward.

\section{Safe recovery}
\begin{description}
\item[Accidentally staged:] After the first commit,
\texttt{git restore --staged FILE} unstages without changing the working copy.
\item[Revert commit:] Starting clean, \texttt{git revert ID} adds an inverse
commit; resolve conflicts if needed.
\item[Need old content:] Inspect \texttt{git show COMMIT\_ID:path/to/file}
first; copy what you need deliberately.
\item[Stash work:] \texttt{git stash push -u} includes untracked files, not
ignored ones. Later use \texttt{git stash apply}, inspect, then drop the
intended stash.
\end{description}
\texttt{git restore FILE} discards uncommitted edits: it is not the same as
unstaging. Hard reset, clean, and force-push can destroy work and are not
part of the beginner lab. Git cannot usually recover edits it never recorded.

\section{Troubleshooting and privacy}
\begin{description}
\item[Push rejected:] Read the error, check branch/remote, and fetch.
Do not immediately force-push.
\item[Branches diverged:] Use \texttt{git fetch origin}, inspect
\texttt{git log --oneline --graph --all}, and agree on a merge/rebase policy.
\item[Identity unknown:] Configure author name/email, not login credentials.
\item[Authentication fails:] Check the credential manager or SSH setup without
displaying secrets on the projector.
\item[Secret committed:] Revoke/rotate immediately. Deleting it in a later
commit does not erase history; follow your organisation's incident procedure.
\end{description}
Keep passwords, tokens, private student details, and confidential research
data out of public repositories. Git alone is not an independent backup.

\section{Further reading}
Read \href{https://git-scm.com/book/en/v2}{Pro Git},
\href{https://git-scm.com/docs}{the Git reference}, and
\href{https://skills.github.com/}{GitHub Skills}.
\end{document}
